\section{Ablation and Analysis}

\subsection{Effect of physical reranking}

The first ablation removes post-generation physical reranking and selects candidates directly from the generative or surrogate stage. This degradation consistently increases the gap between nominal and verified performance, even when the chosen designs appear visually regular. The result confirms that physical reranking is not merely a refinement heuristic; it is a necessary component for aligning the generative model with the true electromagnetic response.

This finding also clarifies why the paper uses the term \emph{physics-consistent}. The generative model alone does not guarantee consistency, particularly in the presence of solution multimodality and sparse sampling of the response manifold. Closed-loop verification is therefore required to transform a generative proposal mechanism into a reliable inverse-design system.

\subsection{Effect of unified response conditioning}

In a second ablation, we replace the full angle--frequency--polarization target with reduced conditioning signals, such as single-slice objectives or compressed scalar descriptors. These variants typically recover part of the desired operator behavior but fail to preserve the coupled structure of the verified response. Angular distortions, spectral drift, or polarization leakage become more pronounced when the target is oversimplified.

This behavior directly supports the response-space thesis of the manuscript. If the objective is compressed too aggressively, the inverse model no longer receives the information required to disambiguate physically distinct solutions. Full response conditioning therefore contributes not just numerical accuracy, but conceptual coherence between the target specification and the underlying device physics.

\subsection{Generative versus deterministic inverse mapping}

A final comparison examines deterministic regressors against conditional generative models. Deterministic baselines often produce reasonable average responses, but they exhibit reduced diversity and weaker best-case performance after verification. By contrast, the generative model provides a candidate ensemble whose members explore multiple regions of the feasible design set, allowing physical reranking to select geometries with improved verified fidelity.

This distinction is especially important for freeform Fourier-operator metasurfaces, where non-uniqueness is a feature of the problem rather than an implementation nuisance. The benefit of generative inversion is therefore not only that it samples more candidates, but that it respects the one-to-many character of the inverse map defined in the unified response space.
